

Classically, there is no useful abelian category of coherent $\mathcal O_X$-modules over a general base.
Synthetically, the category of finitely presented modules over $R$ does not have all kernels,
since $R$ is not coherent \cite[3.2]{topology-draft}, so module bundles with values in finitely presented modules are not abelian.

This section is about finding a suitable replacement category.
This section is joint work of Felix Cherubini, Dan Christensen and Thomas Thorbjørnsen.
The work in this section profited from discussions with Hugo Moeneclaey and Mohamed Barakat.
\rednote{This is work in progress.}

\begin{theorem}
  For an $R$-module $M$ let $M^\ast\colonequiv \Hom_{\Mod{R}}(M,R)$ be its dual.
  Then dualizing a finite co-presentation of $M$:
  \[ M \hookrightarrow R^n \to R^m \]
  yields a finite presentation:
  \[ R^m \to R^n \twoheadrightarrow M^\ast \]
  This induces an anti-equivalence between finitely presented $R$-modules and finitely co-presented $R$-modules. In particular, the canonical map into the double dual, $M\to M^{\ast\ast}$ is an isomorphims for finitely presented and finitely co-presented $R$-modules.
\end{theorem}

\begin{proof}
  \cite{diffgeo-article}
\end{proof}


\begin{lemma}
  \label{presentation-extension}
  Let $M,N$ be finitely presented and $f:M\to N$.
  Then there is an extension of this morphism to all presentations of $M$ and $N$:
  \begin{center}
    \begin{tikzcd}
      R^{m'}\ar[r]\ar[d,dotted,"\exists"] & R^m\ar[->>,r,"\pi_M"]\ar[d,dotted,"\exists"] & M\ar[d, "f"] \\ 
      R^{n'}\ar[r] & R^n\ar[->>,r,"\pi_N"] & N \\ 
    \end{tikzcd}
  \end{center}
  By dualizing, the analogous statement holds for finitely co-presented modules.
\end{lemma}

\begin{proof}
  By linear extension: For a standard basis vector $e_i:R^m$ we merely have $y:R^n$ such that $\pi_N(y)=f(\pi_M(e_i))$. By exactness, we also get $R^{m'}\to R^{n'}$ by linear extension.
\end{proof}

From every commutative square we can extract some exact sequences.
These will be in particular useful to show that finitely presented $R$-modules are closed under cokernels, and finitely copresented $R$-modules are closed under kernels.
We will also use these sequences to show that Adelman modules are closed under tensor products.

\begin{lemma}
  \label{diff-ker-is-ker}
  Suppose we have a commutative square:
  \begin{center}
    \begin{tikzcd}
      A \ar{r}{u} \ar{d}{f} & B \ar{d}{g} \\
      C \ar{r}{v} & D
    \end{tikzcd}
  \end{center}
  The kernel $K$ of the induced map $\mathrm{Ker}(u) \rightarrow \mathrm{Ker}(v)$ is equivalent to the kernel $K$ of $(u, f) : A \to B \oplus C$.
\end{lemma}
\rednote{Is ``difference kernel'' the right name?  I see no subtraction here.  I removed it.}
\rednote{I got rid of intersections, as they create an extra (missing) step, to show that
they are the same as kernels of maps into coproducts.}
\rednote{The difference kernel of f and g is defined as the intersection of their kernels, so the name was more appropriate with the intersections present.}

\begin{proof}
  Observe that $\mathrm{Ker}(u) \subseteq \mathrm{Ker}(gu)$.
  Thus $\mathrm{Ker}(u) = \mathrm{Ker}((u, gu) : A \to B \oplus D)$.
  Therefore we obtain this result by the third isomorphism theorem:
  \begin{center}
    \begin{tikzcd}
      K \ar[equal]{r} \ar[hook]{d} & K' \ar[hook]{d} \\
      \mathrm{Ker}(u) \ar{d} \ar[hook]{r} & A \ar{d}{(u,f)} \ar{r}{(u,gu)} & B \oplus D \ar[equal]{d} \\
      \mathrm{Ker}(v) \ar[hook]{r} & B \oplus C \ar{r}{\mathrm{id} \oplus v} & B \oplus D
    \end{tikzcd}
  \end{center}
  \hfill{\phantom{.}}
\end{proof}

\begin{proposition}
  \label{fcp-kernel-closure}
Finitely copresented modules are closed under kernels.
Dually, finitely presented modules are closed under cokernels.
\end{proposition}

\begin{proof}
  Suppose that $K$ is the kernel of a map $f : M \rightarrow N$, where $M$ and $N$ are finitely copresented.
  By \Cref{presentation-extension}, we can find a linear extension of $f$ to copresentations of $M$ and $N$:
  \begin{center}
    \begin{tikzcd}
      K \ar[hook]{d} \\
      M \ar{d} \ar[hook]{r} & R^{m_1} \ar{r} \ar{d} & R^{m_2} \ar{d} \\
      N \ar[hook]{r} & R^{n_1} \ar{r} & R^{n_2}
    \end{tikzcd}
  \end{center}
  By \Cref{diff-ker-is-ker}, $K$ is also the kernel of the map $R^{m_1} \rightarrow R^{m_2 + n_1}$.
\end{proof}

\begin{lemma}
  \label{kernel-fp-is-cok-fcp}
  Let $f:M\to N$ be a map of finitely presented $R$-modules, then $\ker f$ is the cokernel of a map between finitely copresented $R$-modules.
\end{lemma}

\begin{proof}
  By \Cref{presentation-extension} we can assume $f$ is induced by a square and we construct finitely copresented modules $A,B$ like described below:
  \begin{center}
  \begin{tikzcd}
    A\ar[r,"\pi_1"]\ar[d,"l\times\id"] & R^{m'}\ar[r,"t"]\ar[d,"l"] & R^{n'}\ar[d,"r"]   \\
    B\ar[r,"\pi_1"]\ar[d] & R^m\ar[r,"b"]\ar[d] & R^n\ar[d] \\
    K\ar[r] &M\ar[r,"f"] & N
    \end{tikzcd}
   \end{center}
   $A:=\{(z,y):R^{m'}\times R^{n'}\mid blz=ry\}$, $B:=\{(x,y):R^{m}\times R^{n'}\mid bx=ry\}$.
\end{proof}

We name the following modules after Adelman, since it will turn out that what we construct agrees with the free abelian category on the finite free $R$-modules, and our construction is close to Adelman's general construction of the free abelian category.

The idea is that they might yield an replacement of the algebro-geometric notion of coherent sheaf of modules, which would be a trivial notion over $R$.  

\begin{definition}
  Let $M$ be an $R$-module, then $M$ is an \notion{Adelman-module} if one of the following equivalent statements hold:
  \begin{enumerate}[(i)]
  \item $M$ is merely the kernel of a homomorphism of finitely presented $R$-modules.
  \item $M$ is merely the cokernel of a homomorphism of finitely copresented $R$-modules.
  \item $M$ is the image of a map from a finitely co-presented $R$-module to a finitely presented $R$-module.
  \item $M$ is in the smallest abelian subcategory of $\Mod{R}$ that contains all finite free $R$-modules.
  \end{enumerate}
  We denote the type of Adelman-modules with $\Mod{R}_{\mathrm{ad}}$\index{$\Mod{R}_{\mathrm{ad}}$}, as well as the category of these modules.
\end{definition}
We will show the equivalence of statement (iii) and the first two in ??.
The equivalence with (iv) will follow from \Cref{coherent-R-modules-abelian}.
(iii) will turn out to be self-dual.

\begin{remark}
  Any Adelman-module is wqc (\cite{draft}[Definition 7.1.5]).
\end{remark}

\begin{proof}
  By \cite{draft}[7.1.6, 7.1.13], the category of wqc $R$-modules is abelian and it contains finite free $R$-modules.
\end{proof}

\begin{lemma}
  \label{coherent-double-dual-mono}
  For any Adelman-module $K$, the double dualization map $K\to K^{\ast\ast}$ is a monomorphism.
\end{lemma}

\begin{proof}
  Let $A$ and $B$ be finitely co-presented $R$-modules with a homomorphism $f:A\to B$ such that $K$ is the kernel, then we have the following diagram, which we will explain below:
\begin{center}
\begin{tikzcd}
  B\ar[r,"\sim"] & B^{\ast\ast} \\   
  A\ar[r,"\sim"]\ar[u,"f"] & A^{\ast\ast}\ar[u,"f^{\ast\ast}"] \\   
  K\ar[r,"i"]\ar[u,hook]\ar[ur,hook] & K^{\ast\ast}\ar[u]\ar[l,dotted,bend left=20,"p"] \\   
\end{tikzcd}
\end{center}
$K$ is also the kernel of $f^{\ast\ast}$ by the indicated composition.
By the universal property of this kernel, the dotted map exists and we have $p\circ i=\id$.
\end{proof}

The following lemma will eventually imply that finitely co-presented $R$-modules are projective and, by duality, finitely presented $R$-modules are injective.

\begin{lemma}
  \label{maps-into-coherent-lift}
  Let $A$ and $B$ be finitely co-presented $R$-modules, $K$ be a Adelman-module and $e:A\to K$ a epimorphism.
  Then there is a lift for any homorphism $f:A\to K$.
    \begin{center}
      \begin{tikzcd}
         & C\ar[d,->>,"e"] \\
       A\ar[r,"f"]\ar[ru,dotted,"\exists"] & K
      \end{tikzcd}
    \end{center}
\end{lemma}
\rednote{I assume that $B$ should be $C$.  Is it correct that we need both $A$ and $C$ to be finitely co-presented?}

\begin{proof}
  We will prove the lemma using the following steps:
  \begin{enumerate}[(i)]
  \item \label{coherent-fcp-lift-existence} There is a lift on the level of the underlying sets, i.e.\ a map $\tilde{f}$ as indicated in the commutative diagram, which is not necessarily $R$-linear:
    \begin{center}
      \begin{tikzcd}
         & C\ar[d,->>,"e"] \\
       A\ar[r,"f"]\ar[ru,dotted,"\exists \tilde{f}"] & K 
      \end{tikzcd}
    \end{center}
  \item Any lift $\tilde{f}$ from \ref{coherent-fcp-lift-existence} can be decomposed into a $R$-linear part $f_L:A\to C$, and a non-linear part $f_{NL}\coloneq\tilde{f}-f_L$.
  \item $e$ maps $f_{NL}$ to $0$, so $f_L$ is a $R$-linear lift.
  \end{enumerate}
  We will now prove these statements.
  \begin{enumerate}[(i)]
  \item Let $M$ be the kernel of $e$, then $M$ is wqc.
  By \cite{draft}[Theorem 7.3.6] we have $H^1(A,M)=0$, since $A$, like any finitely co-presented $R$-module, is an affine scheme.

  By definition, we have $H^1(A,M)=||A\to K(M,1)||_0$ and therefore that $A\to K(M,1)$ is connected.
  For $x:K$ let $E_x$ denote the fiber of $e$ over $x$, then $E_x$ is a $M$-torsor or in other words an element of $K(M,1)$.
  So $E\circ f:A\to K(M,1)$ and we have $|| E\circ f = (\_\mapsto M) ||$ by the connectedness noted above.
  This means that the dependent type $E\circ f$ merely has a section, since $(\_\mapsto M)$ has a section.
  This proves existence of the dotted map.
  \item By \cite{diffgeo-article}, we know that the differential of $\tilde{f}$ is such a linear map,
  but a direct computation is also possible, which we will carry out in the following.
  
  Let $A$ and $C$ be given by $A=\Spec R[X_1,\dots,X_n]/(P_1,\dots,P_r)$ and $B=\Spec R[Y_1,\dots,Y_m]/(Q_1,\dots,Q_s)$.
  $\tilde{f}$ is represented by a homomorphism $\varphi$ between these $R$-algebras.
  Choosing lifts of the finitely many $\varphi(Y_i)$ we get polynomials $N_1,\dots,N_m$
  such that for all $j$ we have $Q_j(N_1,\dots,N_m)\in (P_1,\dots,P_r)$.
  Let $L_1,\dots,L_m$ be given by the linear parts of $N_1,\dots,N_m$.
  Then by analysing the degrees of the involved polynomials, we have $Q_j(L_1,\dots,L_m)\in (P_1,\dots,P_r)$, so the $L_1,\dots,L_m$ induce a homorphism
  \[
  \varphi_L: R[Y_1,\dots,Y_m]/(Q_1,\dots,Q_s) \to R[X_1,\dots,X_n]/(P_1,\dots,P_r)
  \]
  which induces a linear map $f_L:A\to B$.
  Let us note that the difference $f_{NL}\coloneq \tilde{f}-f_L$ is given by monomials $S_i$ which are all of degree at least two.
  \item We know that $e\circ \tilde{f}$ is a linear map, so $e\circ f_{NL}$ is linear as well.
  Let $f_{NL}=\sum_i S_i$ be a decomposition into monomials with degrees $n_i\geq 2$ as above.
  For all $t:R$ and $x:A$ we have
  \[
  \sum_i te(S_i(x))=te(\sum_i S_i(x))=e(\sum_i S_i(tx))=e(\sum_i t^{n_i}S_i(x))=\sum_i t^{n_i}e( S_i(x))
  \]
  So by taking the difference we have a polynomial $P(t)\coloneq \sum_i te(S_i(x))-\sum_i t^{n_i}e( S_i(x))$ in $t$ with coefficients in $K$, such that $P(t)=0$ for all $t$.
  Since all $e(S_i(x))$ appear as coefficients of $P$, we can conclude $e\circ f_{NL}=0$, from $P$ being zero.

  Let $f:K\to R$ be an arbitrary linear map, then $f(P(t)):R\to R$ has to be the zero polynomial,
  so $f$ maps all $e(S_i(x))$ to zero.
  This means that the double dualization $K\to K^{\ast\ast}$ maps each $e(S_i(x))$ to zero,
  but by \Cref{coherent-double-dual-mono} it is injective, so all $e(S_i(x))$ are zero.
  \end{enumerate}
\end{proof}

\begin{lemma}
  \label{coherent-kernels}
  The cokernel of a homomorphism between Adelman-modules is coherent.  
\end{lemma}

\begin{proof}
  Let $K,L$ be Adelman-modules and $f:K\to L$ a homomorphism.
  Choose finitely co-presented $R$-modules $A,B,C,D$ and homomorphisms to present $K$ and $L$ as cokernles.
  By the preceeding \Cref{maps-into-coherent-lift} there are lifts yielding this diagram:
  \begin{center}
    \begin{tikzcd}
      A\ar[d]\ar[r,dotted] & C\ar[d,"r"] & \\
      B\ar[d,->>]\ar[r,dotted,"b",swap] & D\ar[d,->>]\ar[rd,->>,"g"] & \\
      K\ar[r,"f",swap] & L\ar[r,->>,"\mathrm{cok}",swap] & \mathrm{Cok}
    \end{tikzcd}
  \end{center}
  Now we conclude by showing that the composition $g$ is the cokernel of a homorphism of the finitely co-presented $R$-modules $B\oplus C$ and $D$.
  By diagram-chasing, $b\oplus r:B\oplus C\to D$ surjects onto the kernel of $g$.
\end{proof}

\begin{lemma}
  \label{fp-kernel-dual-surjective}
  Let $M,N$ be finitely presented $R$-modules, $f:M\to N$ a homomorphism and $K$ its kernel,
  then the dual of the canonical map $K\to M$ is surjective.
\end{lemma}

\begin{proof}
  It is possible \rednote{TODO: add construction} to construct a diagram of the following shape, where the two left columns and all rows are exact:
  \begin{center}
    \begin{tikzcd}
      A\ar[r,hook]\ar[d] & R^{m_1}\ar[r,"t"]\ar[d,"l"] & R^{n}\ar[d,"\id"] \\
      B\ar[r,hook]\ar[d,->>] & R^{m_2}\ar[r,"b"]\ar[d,->>] & R^{n} \\
      K\ar[r,hook] & M\ar[r] & N
    \end{tikzcd}
  \end{center}
  Then a linear map $L:K\to R$ restricts to a map $L_1:B\to R$ which is zero on $A$.
  By \cite{diffgeo-article}[Lemma 2.3.7] it is possible to extend $L_1$ to $L_2:R^{m_2}\to R$.
  $L_2$ will descend to $M$, if and only if it vanishes on the image of $r$.
  This will generally not be the case, but we can construct a linear correction $c:R^{m_2}\to R$ such that $L_2-c$ has this property.

  To construct $c$, let $c_0\coloneq L_2\circ l$, note that $c_0$ vanishes on $A$ and therefore extends to the image of $t$, to a linear map $c_1:\mathrm{im}(t)\to R$.
  By \cite{diffgeo-article}[Lemma 2.3.8], $c_1$ extends to a map $c_2:R^{n}\to R$.
  Now $c\coloneq c_2\circ b$ is a linear map on $R^{m_2}$, that is zero on $B$,
  so $L_2-c$ is still an extension of $L_1$ and $L_2-c$ vanishes on $A$.
\end{proof}

\begin{theorem}
  \label{coherent-self-dual}
  Let $K$ be a Adelman-module, then the canonical map $K\to K^{\ast\ast}$ is an isomorphism.
  In particular, dualization is an anti-equivalence of the category of Adelman-modules.
\end{theorem}

\begin{proof}
  We can assume $K$ is in a sequence of $R$-modules:
  \begin{center}
    \begin{tikzcd}
      A\ar[r,"f"] & B\ar[r,->>] & K
    \end{tikzcd}
  \end{center}
  where $A,B$ are finitely co-presented and $B\to K$ is the cokernel of $f$.
  Taking double duals and using \Cref{fp-kernel-dual-surjective} and \Cref{coherent-double-dual-mono} we end up in this situation:
  \begin{center}
    \begin{tikzcd}
      A\ar[r,"f"] & B\ar[r,->>]\ar[d,"\sim"] & K\ar[d,hook] \\
      & B^{\ast\ast} \ar[r,->>] & K^{\ast\ast} 
    \end{tikzcd}
  \end{center}
  So $K\to K^{\ast\ast}$ is an isomorphism.
\end{proof}

\begin{theorem}
  \label{coherent-R-modules-abelian}
  The category of Adelman-modules is abelian.
\end{theorem}

\begin{proof}
  By \Cref{coherent-kernels} Adelman-modules are closed under kernels.
  By the duality from \Cref{coherent-self-dual}, this also shows the closure under cokernels.
\end{proof}

\begin{lemma}
  \label{fcp-projective-fp-injective}
  Any finitely co-presented $R$-module $A$ is projective in the category of of Adelman-modules:
  For all surjections $K\to K'$ with $K,K':\Mod{R}_{\mathrm{ad}}$ and homomorphisms $f:A\to K'$ there exists a lift:
  \begin{center}
  \begin{tikzcd}
    & K\ar[d,->>] \\
    A\ar[r,"f",swap]\ar[ru,dotted,"\exists"] & K'
\end{tikzcd}
\end{center}
  Dually, any finitely presented $R$-module is injective in $\Mod{R}_{\mathrm{ad}}$.
\end{lemma}

\begin{proof}
   Use \Cref{maps-into-coherent-lift} and that for any Adelman-module like $K$, there is a finitely co-presented $R$-module surjecting onto it.
\end{proof}

\begin{remark}
  By \cite[Theorem 1.15]{adelman-construction}, \Cref{fcp-projective-fp-injective} implies that the Adelman-modules are the free abelian category over the finite free $R$-modules.
  This means the Adelman-modules have the following universal property:
  For every abelian category $\mathcal{A}$ and all additive functors $F:\Mod{R}_{\mathrm{ff}}\to \mathcal{A}$, there exists a unique exact functor $K$ like follows:
  \begin{center}
  \begin{tikzcd}
    \Mod{R}_{\mathrm{ff}}\ar[r]\ar[rd,"F",swap] & \Mod{R}_{\mathrm{ad}}\ar[d,dotted,"K"] \\
    & \mathcal{A}
\end{tikzcd}
\end{center}
\end{remark}

\begin{remark}
  \label{hom-closure}
  For Adelman-modules $M,N$, the $R$-module $\Hom_{\Mod{R}}(M,N)$ is an Adelman-module as well.
\end{remark}

\begin{proof}
  The proof in \cite{lombardi-quitte}[Chapter IV, 4.12] works, if ``coherent'' is replaced by ``Adelman''.
  
  First we use their proof to show that for finitely presented $R$-modules $M,N$ the $R$-module $\Hom_{\Mod{R}}(M,N)$ is Adelman. Morphisms $f:M\to N$ are presented by squares:
  \begin{center}
  \begin{tikzcd}
        R^m \ar[r]\ar[d,"\varphi"] & R^{m'}\ar[d,"\psi"] \\
        R^n\ar[r] & R^{n'}
\end{tikzcd}
\end{center}
  So we have a finite free module of pairs of morphisms $(\varphi,\psi)$ and a submodule of pairs such that the square above commutes.
  This submodule $S$ is the kernel of a linear map and therefore Adelman.
  We have a surjection $\pi:S\to \Hom_{\Mod{R}}(M,N)$, so $\Hom_{\Mod{R}}(M,N)$ is Adelman,
  if $\pi$ is a cokernel of a map of Adelman-modules.
  This is the case, since there is a surjection onto $\ker{\pi}$ from the finite free $R$-module of linear maps which splitting the square:
  \begin{center}
  \begin{tikzcd}
        R^m \ar[r]\ar[d,"\varphi"] & R^{m'}\ar[d,"\psi"]\ar[ld,"s"] \\
        R^n\ar[r] & R^{n'}
\end{tikzcd}
\end{center}
  Then we can reuse the same argument to show the statement of the remark.
\end{proof}

\begin{lemma}
  \label{H1-on-adelman-module}
  \rednote{Beweis leider Quatsch}
  Let $K:\Mod{R}_{\mathrm{ad}}$ and $M:K\to \Mod{R}_\mathrm{wqc}$ then
  \[
    H^1(K,M)=0
  \]
\end{lemma}

\begin{proof}
  Let first $K:\Mod{R}$ be in an exact sequence:
  \begin{center}
    \begin{tikzcd}
      F\ar[r,"e"] & K\ar[r]\ar[l,bend left,"s",dotted] & 0
    \end{tikzcd}
  \end{center}
  with $F:\Mod{R}_\mathrm{fcp}$ and $M:K\to\Mod{R}_{\mathrm{wqc}}$ and $s$ exists because $K$ is Adelman using \Cref{maps-into-coherent-lift}.
  Then we have $H^1(F,M\circ e)=0$, meaning that for any $\chi:(x:K)\to K(M_x,1)$ there exists $t:(y:F)\to \chi_{e(y)}=\ast$. Then $t\circ s:(x:K)\to \chi_x=\ast$, so $H^1(K,M)=0$.
\end{proof}

\begin{proposition}
  \label{fp-extension-adelman}
  Let
  \[
    0\to M \to N \to L \to 0
  \]
  be an exact sequence such that $M :\Mod{R}_{\mathrm{ad}}$, and $L : \Mod{R}_{\mathrm{fp}}$.
  Then $N$ is Adelman.
\end{proposition}

\begin{proof}
  This proof is parallel to the horseshoe lemma.
  Let $R^{n_2} \rightarrow R^{n_1} \twoheadrightarrow L$ be a presentation of $L$.
  Let also $P_2 \rightarrow P_1 \twoheadrightarrow M$ be a presentation of $M$ by fcp-modules.
  Because $R^{n_1}$ is projective, we can lift the morphism $R^{n_1} \rightarrow L$ along $N \twoheadrightarrow L$.
  This yields a map of short exact sequences.
  \begin{center}
    \begin{tikzcd}
      P_1 \ar[hook]{r} \ar[two heads]{d} & P_1 \oplus R^{n_1} \ar[two heads]{r} \ar[dashed, two heads]{d} \ar[phantom]{rd}[very near end]{\ulcorner} & R^{n_1} \ar[two heads]{d} \ar[dashed]{ld} \\
      M \ar[hook]{r} & N \ar[two heads]{r} & L
    \end{tikzcd}
  \end{center}
  One can show that the map $P_1 \oplus R^{n_1} \rightarrow N$ is epi with a diagram chase.
  The lower rightmost square is also a pushout, since the maps on the kernels is epic.
  Therefore taking kernels of the vertical epimorphisms yields another exact sequence, and we repeat the above argument to conclude.
  \begin{center}
    \begin{tikzcd}
      P_2 \ar[hook]{r} \ar[two heads]{d} & P_2 \oplus R^{n_2} \ar[two heads]{r} \ar[dashed, two heads]{d} & R^{n_2} \ar[dashed]{ld} \ar[two heads]{d} \\
      K_M \ar[hook]{r} & K_N \ar[two heads]{r} & K_L
    \end{tikzcd}
  \end{center}
  Thus we get a presentation $P_2 \oplus R^{n_2} \rightarrow P_1 \oplus R^{n_1} \twoheadrightarrow N$ of $N$ by fcp-modules.
\end{proof}

\begin{corollary}
  Every short exact sequence
  \begin{center}
    \begin{tikzcd}[cramped]
      I_1 \ar[hook]{r} & E \ar[two heads]{r} & I_2
    \end{tikzcd}
  \end{center}
  with $I_i$ fp-modules is split.
  Consequently $\textup{Ext}^1_R(I_2,I_1) = 0$.
\end{corollary}

\begin{proof}
  From \Cref{fp-extension-adelman}, $E$ is an Adelman-module.
  By \Cref{fcp-projective-fp-injective}, the inclusion $I_1 \rightarrow E$ splits, so $E = I_1 \oplus I_2$.
\end{proof}

\begin{lemma}
  \label{adelman-ext-pb-red}
  Suppose we have a commutative diagram
  \begin{center}
    \begin{tikzcd}
      B \ar[equal]{d} \ar[hook]{r} & p^*E \ar[two heads]{r} \ar[two heads]{d} \ar[phantom]{rd}[very near start]{\lrcorner}[very near end]{\ulcorner} & P \ar[two heads]{d}{p} \\
      B \ar[hook]{r} & E \ar[two heads]{r} & A
    \end{tikzcd}
  \end{center}
  in which $B$, $A$ and $P$ are Adelman modules.
  Then $E$ is Adelman if and only if $p^*E$ is so.
\end{lemma}

\begin{proof}
  Suppose that $E$ is Adelman.
  Since $p^*E$ is a pullback of Adelman modules, it is itself Adelman.

  Now assume that $p^*E$ is Adelman.
  By the third isomorphism theorem we can prove that $E$ is Adelman.
  In the following diagram, the objects $0^*_Ap^*E$ is the kernel of $p^*E \rightarrow E$, $0^*_AE$ is the kernel of $E \rightarrow A$, and $0^*_Ep^*E$ is the kernel of $p^*E \rightarrow E$.
  \begin{center}
    \begin{tikzcd}
      0^*_Ep^*E \ar[equal]{r} \ar[hook]{d} & 0^*_Ep^*E \ar[hook]{d} \\
      0^*_Ap^*E \ar[hook]{r} \ar[two heads]{d} & p^*E \ar[two heads]{r} \ar[two heads]{d} & A \ar[equal]{d} \\
      B \ar[hook]{r} & E \ar[two heads]{r} & A
    \end{tikzcd}
  \end{center}
  Chasing around this diagram we can see that $0^*_Ap^*E$ is Adelman, because $p^*E$ is.
  Furthermore, $0^*_Ep^*E$ is Adelman, because $0^*_Ap^*E$ is Adelman.
  It then follows that $E$ is Adelman, because it is the cokernel of two Adelman modules.
\end{proof}

In fact, the same proof shows:

\begin{lemma}
  \label{adelman-ext-quotient-adelman}
  Suppose that $E$ is an extension of Adelman modules, and that there is an
  Adelman module $P$ and an epimorphism $P \to E$.  Then $E$ is Adelman.
\end{lemma}

\begin{proof}
Follow the same proof, with $p^*E$ replaced by $P$.
\end{proof}

\begin{proposition}
  \label{adelman-ext-problem-red}
  Adelman modules are closed under extension if and only if for any fcp-module $P$ and fp-module $I$, $\textup{Ext}^1_R(P,I) = 0$.
\end{proposition}

\begin{proof}
  Assume that Adelman modules are closed under extensions.
  Then it is clear that any extension by $I$ and $P$ is split, as $I$ is injective among the Adelman modules \Cref{fcp-projective-fp-injective}.

  Assume now that $\textup{Ext}^1(P,I) = 0$ for any such modules $P$ and $I$.
  Consider an extension problem
  \begin{center}
    \begin{tikzcd}
      M \ar[hook]{r} & N \ar[two heads]{r} & L
    \end{tikzcd}
  \end{center}
  where $M$ is Adelman and is copresented by $M \hookrightarrow I_1 \rightarrow I_2$, and $L$ is Adelman and is presented by $P_2 \rightarrow P_1 \twoheadrightarrow L$.
  By \Cref{adelman-ext-pb-red}, we can take the pullback of the extension along $p : P_1 \twoheadrightarrow L$, we get another associated extension.
  \begin{center}
    \begin{tikzcd}
      M \ar[equal]{d} \ar[hook]{r} & p^*N \ar[two heads]{d} \ar[two heads]{r} \ar[phantom]{rd}[very near start]{\lrcorner} & P_1 \ar[two heads]{d} \\
      M \ar[hook]{r} & N \ar[two heads]{r} & L
    \end{tikzcd}
  \end{center}
  Now $N$ is Adelman if and only if $p^*N$ is so.

  Dually, we can take a pushout of $M \hookrightarrow p^*N \twoheadrightarrow P_1$  to obtain the diagram.
  \begin{center}
    \begin{tikzcd}
      M \ar[hook]{r} \ar[hook]{d} \ar[phantom]{rd}[very near end]{\ulcorner} & p^*N \ar[two heads]{r} \ar[hook]{d} & P_1 \ar[equal]{d} \\
      I_1 \ar[hook]{r} & i_*p^*N \ar[two heads]{r} & P_1
    \end{tikzcd}
  \end{center}
  In this case $p^*N$ is Adelman if and only if $i_*p^*N$ is Adelman.

  Now, since we assumed that $\textup{Ext}^1_R(P_1,I_1) = 0$, it is necessary that $i_*p^*N = I_1 \oplus P_1$, which is Adelman.
  Thus, $N$ is Adelman as desired.
\end{proof}

\rednote{We do not know how to compute these extension groups yet.}

An idea that might be useful is to apply dimension shift. Let
\begin{center}
  \begin{tikzcd}
    P \ar[hook]{r} & R^{n_2} \ar{r} & R^{n_1} \ar[two heads]{r} & I_P
  \end{tikzcd}
\end{center}
be a two-step projective resolution of $I_P$.
By dimension shift, we get an isomorphism $\textup{Ext}^1(P,M) = \textup{Ext}^3(I_P,M)$ for any $R$-module $M$.

Using \Cref{adelman-ext-pb-red}, we can deduce a stronger result than \Cref{fp-extension-adelman}.

\begin{proposition}
  \label{fg-extension-adelman}
  Let
  \begin{center}
    \begin{tikzcd}
      0 \ar{r} & M \ar[hook]{r} & N \ar[two heads]{r} & L \ar{r} & 0
    \end{tikzcd}
  \end{center}
  be an exact sequence such that $M,L : \Mod{R}_{\textup{ad}}$ and $L$ is fg.
  Then $N$ is Adelman.
\end{proposition}

\begin{proof}
  Consider the pullback
  \begin{center}
    \begin{tikzcd}
      M \ar[equal]{d} \ar[hook]{r} & p^*N \ar[two heads]{d} \ar[two heads]{r} & R^n \ar[two heads]{d} \\
       M \ar[hook]{r} & N \ar[two heads]{r} & L
    \end{tikzcd}
  \end{center}
  Since the epi $p^*N \twoheadrightarrow R^n$ splits, $p^*N$ is an Adelman module.
  By \Cref{adelman-ext-pb-red}, $N$ is Adelman as well.
\end{proof}

The following method is an idea to extend the current results beyond finitely generated Adelman modules to countably generated ones.
In order to make this work we will assume the axiom of dependent choice, but it is not known if we can get around this.
Suppose that $M$ is a countably generated Adelman module, and let $g : R[\mathbb{N}] \twoheadrightarrow M$ be the generating map.
The following sequence of submodules $M_n := \langle g(0), \ldots, g(n) \rangle$ with the canonical injections $i_n : M_n \hookrightarrow M_{n+1}$ has $M$ as its colimit.
We can then express the colimit by the following short exact sequence
\begin{center}
  \begin{tikzcd}
    \underset{n : \mathbb{N}}{\bigoplus} M_n \ar[hook]{r}{\nabla} & \underset{n : \mathbb{N}}{\bigoplus} M_n \ar[two heads]{r} & M
  \end{tikzcd}
\end{center}
where $\nabla(n,m) = i_n(m) - m$.

\begin{proposition}
  Assume that dependent choice holds.
  Let $M$ be a countably generated Adelman module.
  Then for any finitely presented module $I$, $\textup{Ext}^1(M,I)=0$.
\end{proposition}

\begin{proof}
  Applying $\textup{Hom}(\cdot, I)$ to the above short exact sequence yields a long exact sequence
  \begin{center}
    \begin{tikzcd}
      \textup{Hom}(M,I) \ar[hook]{r} & \underset{n : \mathbb{N}}{\prod} \textup{Hom}(M_n, I) \ar{r}{\Delta} & \underset{n : \mathbb{N}}{\prod} \textup{Hom}(M_n, I) \ar[looseness = 1, out = -10, in = 170]{dll} \\
      \textup{Ext}^1(M,I) \ar{r} & \underset{n : \mathbb{N}}{\prod} \textup{Ext}^1(M_n, I) \ar{r} & \underset{n : \mathbb{N}}{\prod} \textup{Ext}^1(M_n, I)
    \end{tikzcd}
  \end{center}
  where $\Delta = \nabla^*$.
  By \Cref{fg-extension-adelman}, we know that $\textup{Ext}^1(M_n, I) = 0$ for all $n : \mathbb{N}$.
  Thus, $\textup{Ext}^1(M,I) = 0$ if and only if $\Delta$ is epic.

  We prove that $\Delta$ is epic using dependent choice.
  Let $\Delta_k : \underset{0 \leq i \leq k + 1}{\prod} \textup{Hom}(M_i, I) \rightarrow \underset{0 \leq i \leq k}{\prod} \textup{Hom}(M_i, I)$ be the restriction of $\Delta$ to the indicated factors, and set $\Delta_{-1}$ to be the terminal map.
  \rednote{I think the subscripts on the products need to be the same, since one of the components of $\nabla$ is $-\textup{id}$.  Unless you mean to co-restrict in the codomain as well?}
  First we show that each $\Delta_k$ is epi.
  Let $g : \underset{0 \leq i \leq k}{\prod} \textup{Hom}(M_i, I)$ be a sequence of homomorphisms.
  The map $\textup{Hom}(M_{k+1},I) \rightarrow \textup{Hom}(M_k, I)$ is epi, because $I$ is injective among the Adelman modules, using \Cref{fcp-projective-fp-injective}.
  For our lift, let $f_0 : M_0 \rightarrow I$ be some map, say the $0$-map.
  Then we inductively make each lift $f_{i+1}$ by lifting $g_i-f_i$ for each $i \leq k$.
  \rednote{I believe that since this induction is finite, we can commute out the pi of the propositional truncation and proceed with the induction as I described.}

  We now need to show that for any $k$, the map into the following pullback
  \begin{center}
    \begin{tikzcd}
      \underset{0 \leq i \leq k + 2}{\prod}\textup{Hom}(M_i, I) \ar[two heads, bend right, looseness = 0.5]{ddr} \ar[two heads,bend left, looseness = 1, out = 15]{rrd}{\Delta_{k+1}} \ar[dashed]{rd} \\
      & \textup{PB} \ar[two heads]{d} \ar[two heads]{r} \ar[phantom]{rd}[very near start]{\lrcorner} & \underset{0 \leq i \leq k + 1}{\prod} \textup{Hom}(M_i, I) \ar[two heads]{d} \\
      & \underset{0 \leq i \leq k + 1}{\prod} \textup{Hom}(M_i, I) \ar[two heads]{r} & \underset{0 \leq i \leq k}{\prod} \textup{Hom}(M_i, I)
    \end{tikzcd}
  \end{center}
  is epi.
  \rednote{I don't understand what this pullback has to do with our goal.}
  An element of the pullback is given by two sequences $(f,g)$ such that $\Delta_k(f)_i = g_i$ for all $i \leq k$.
  Thus the sequence $f$ determines a valid lift, only missing the $(k+1)$-st component.
  We set $f_{k+1}$ to be a lift of $f_k - g_k$ along $i_k$.
  This puts us in a position where we can use dependent choice, so $\Delta$ is epic.

  \rednote{Can't we just show directly that $\Delta$ is epic using the method in the second-last paragraph?
  We start with $g : \prod_i \textup{Hom}(M_i, I)$, and define $f$ by making a sequence of choices?}
\end{proof}

\begin{lemma}
  \rednote{TODO: This is just an idea}.
  Let
  \[
    0\to M \to N \to L \to 0
  \]
  be an exact sequence and $M,L:\Mod{R}_{\mathrm{ad}}$.
  Then $N$ is Adelman.
\end{lemma}

\begin{proof}
  By \Cref{H1-on-adelman-module} we have $H^1(L,M)=0$ and therefore a not necessarily linear section $s:L\to N$.
\end{proof}

\rednote{The following can be written out in more detail. Here the Tor-groups are defined purely in terms of resolutions. All results about exact sequences should be recovered from the horseshoe lemma and the snake lemma.}
Adelman modules are also closed under tensor products.
Let $M$ and $N$ be two Adelman modules that are presented by the finitely copresented modules $P_1, P_2$ and $Q_1, Q_2$, respectively.
Because the tensor product is right exact, it preserves presentations.
Consider the following commutative square
\begin{center}
  \begin{tikzcd}
    P_2 \otimes Q_2 \ar{r} \ar{d} & P_2 \otimes Q_1 \ar{d} \\
    P_1 \otimes Q_2 \ar{r} & P_1 \otimes Q_1
  \end{tikzcd}
\end{center}
By the dual of \Cref{diff-ker-is-ker}, we obtain the following exact sequence
\begin{center}
  \begin{tikzcd}
    (P_1 \otimes Q_2) \oplus (P_2 \otimes Q_1) \ar{r} & P_1 \otimes Q_1 \ar[two heads]{r} & M \otimes N
  \end{tikzcd}
\end{center}
If the first two terms of this sequence are Adelman, then so is the last one.
Thus it suffices to show that the tensor of finitely copresented modules is Adelman to show that all Adelman modules are closed under tensor products.
We will prove something slightly stronger, namely that the tensor of fcp modules are themselves fcp.

\begin{lemma}
  Let $P$ and $Q$ be fcp modules.
  Then $P \otimes Q$ is fcp as well.
\end{lemma}

\begin{proof}
  Let $Q \hookrightarrow R^{n_1} \rightarrow R^{n_2}$ be a copresentation of $Q$.
  It suffices to show that the map $P \otimes Q \rightarrow P \otimes R^{n_1}$ is monic, because $P \otimes R^{n_i}$ is fcp for $i = 0$ and $1$.
  (This uses that fcp modules are closed under kernels.)

  Let $M$ be the image of the map $R^{n_1} \rightarrow R^{n_2}$.
  Since both $Q$ and $R^{n_1}$ are projective in the Adelman category, the short exact sequence
  \begin{center}
    \begin{tikzcd}
      Q \ar[hook]{r} & R^{n_1} \ar[two heads]{r} & M
    \end{tikzcd}
  \end{center}
  is a projective resolution.
  Since the tensor product is right exact we get an exact sequence after applying $P \otimes (\cdot)$
  \begin{center}
    \begin{tikzcd}
      \textup{Tor}_1(P,M) \ar[hook]{r} & P \otimes Q \ar{r} & P \otimes R^{n_1} \ar[two heads]{r} & P \otimes M
    \end{tikzcd}
  \end{center}
  As $\textup{Tor}_1(\cdot, \cdot)$ is balanced, we can compute it starting from $P$, but then we see that $\textup{Tor}_1(P,M) = 0$, so the map $P \otimes Q \rightarrow P \otimes R^{n_1}$ is monic.
\end{proof}

\subsection{Strange modules}\label{strange-modules}

\rednote{These are meant to be ideas/WIP that might not appear in a final draft.}

Given a type $U$, one can form the following modules:
\begin{equation}\label{eq:strange-sup}
  \begin{tikzcd}
    0 \ar[r] & K^U \ar[r,hook] & R \ar[rr,"c"] \ar[rd,->>] && R^U \\
             &                  &            & R/K^U \ar[ru]
  \end{tikzcd}
\end{equation}
where $c$ sends $r$ to the constant function with value $r$,
$K^U$ is defined to be the kernel of $c$, and $R/K^U$ is the image of $c$.
Since $c$ is a ring homomorphism, $K^U$ is an ideal, and $R/K^U$ is also a ring.

We can also form
\begin{equation}\label{eq:strange-sub}
  \begin{tikzcd}
    0 \ar[r] & R_U \ar[r,hook] & R \ar[r,->>] & R/R_U \ar[r] & 0
  \end{tikzcd}
\end{equation}
where $R_U$ is the submodule defined by the predicate sending
$x : R$ to $(x = 0) \lor \propTrunc{U}$.
For this sequence, we may as well assume that $U$ is a proposition.
Again, I think $R_U$ is an ideal and the quotient is a ring.

In total, this gives five modules that depend on a type $U$,
three of which are rings.

There is a relationship between \eqref{eq:strange-sub} for $\lnot U$
(top row in the following diagram)
and \eqref{eq:strange-sup} for $U$ (bottom row):
\begin{equation}\label{eq:strange-compare}
  \begin{tikzcd}
    0 \ar[r] & R_{\lnot U} := \sum_{x : R} (x = 0) \lor \lnot U \ar[r,hook] \ar[d,hook]
    & R \ar[r,->>] \ar[d,equal] & R/R_{\lnot U} \ar[r] \ar[d,->>] & 0 \\
    0 \ar[r] & K^U \ar[r,hook]
    & R \ar[r,->>] & R/K^U \ar[r] & 0
  \end{tikzcd}
\end{equation}
The downwards inclusion on the left is easy to check:  Let $x : R$.  If $x = 0$,
then $x$ is clearly in the submodule $K^U$, and if $\lnot U$, then $R^U = 0$, so $K^U$ is all
of $R$, so $x$ is in $K^U$.
This map is a monomorphism, since it factors one.
The epimorphism on the right follows formally.

Specializing \eqref{eq:strange-sub} to $U = V(r) = \Spec R/(r) = (r = 0)$,
for a chosen $r : R$, gives the top row of this diagram,
which is related to the kernel-image exact sequence coming from the
map $\varphi_r : R \to R$ that multiplies by $r$:
\begin{equation*}
  \begin{tikzcd}
    0 \ar[r] & R_{V(r)} := \sum_{x : R} (x = 0) \lor (r = 0) \ar[r,hook] \ar[d,hook]
    & R \ar[r,->>] \ar[d,equal] & R/R_{V(r)} \ar[r] \ar[d,->>] & 0 \\
    0 \ar[r] & \ker \varphi_r := \sum_{x : R} (xr = 0) \ar[r,hook]
    & R \ar[r,->>] & \im \varphi_r \ar[r] & 0
  \end{tikzcd}
\end{equation*}
The downwards inclusion on the left is easy to see, and the downwards
epimorphism on the right follows formally.
(This is not a special case of \eqref{eq:strange-compare}.)

Recall that $D(r) = \lnot V(r) = (r \neq 0) = \inv(r) = \Spec R[1/r]$, so
\eqref{eq:strange-compare} with $U = V(r)$ involves $V(r)$ and $D(r)$:
\begin{equation}\label{eq:compare-Vr}
  \begin{tikzcd}
    0 \ar[r] & R_{D(r)} := \sum_{x : R} (x = 0) \lor (r \neq 0) \ar[r,hook] \ar[d,hook]
    & R \ar[r,->>,"q"] \ar[d,equal] & R/R_{D(r)} \ar[r] \ar[d,->>] & 0 \\
    0 \ar[r] & K^{V(r)} := (r) \ar[r,hook]
    & R \ar[r,->>,"c"] & R^{V(r)} = R/(r) \ar[r] & 0
  \end{tikzcd}
\end{equation}
Here note that $R^{V(r)} = R/(r)$ since $V(r) = \Spec R/(r)$, and so the kernel of $c$ is $(r)$.
Also, the map $c$ is epi in this case, so $R^{V(r)}= R/K^{V(r)}$.

The map $q$ in the previous diagram has the property that $\lnot (q(r) \neq 0)$.
Indeed, assume that $q(r) \neq 0$.  Then $r \neq 0$, and therefore $R/R_{D(r))} = 0$,
so $q(r) = 0$, a contradiction.
\textbf{Questions:}
Does this mean that $q(r)$ is nilpotent in the ring $R/R_{D(r)}$?
Or does the relationship between not non-zero and nilpotence only work for the ring $R$?
Is $q$ the universal map for which $\lnot (q(r) \neq 0)$?

Next we can consider \eqref{eq:strange-compare} with $U := D(r)$, so
$\lnot U = \lnot (r \neq 0)$.  We get
\begin{equation*}
  \begin{tikzcd}
    0 \ar[r] & R_{\lnot D(r)} := \sum_{x : R} (x = 0) \lor \lnot (r \neq 0) \ar[r,hook] \ar[d,hook]
    & R \ar[r,->>] \ar[d,equal] & R/R_{\lnot D(r)} \ar[r] \ar[d,->>] & 0 \\
    0 \ar[r] & K^{D(r)} = \sum_{x : R} \exists_n (r^n x = 0)  \ar[r,hook]
    & R \ar[r,->>] & R/K^{D(r)} \ar[r] & 0
  \end{tikzcd}
\end{equation*}
where $R/K^{D(r)}$ is the image of $R$ in $R^{D(r)} = R[1/r]$.
The displayed description of $K^{D(r)}$ follows from the construction of $R[1/r]$
using equivalence classes of fractions.  \rednote{Citation?}

Finally, consider \eqref{eq:strange-compare} with $U := \lnot D(r)$, noting that
$\lnot \lnot D(r) = D(r)$ (since $D(r)$ is already a negation):
\begin{equation*}
  \begin{tikzcd}
    0 \ar[r] & R_{D(r)} := \sum_{x : R} (x = 0) \lor (r \neq 0) \ar[r,hook] \ar[d,hook]
    & R \ar[r,->>,"q"] \ar[d,equal] & R/R_{D(r)} \ar[r] \ar[d,->>] & 0 \\
    0 \ar[r] & K^{\lnot D(r)} \ar[r,hook]
    & R \ar[r,->>,"q'"] & R/K^{\lnot D(r)} \ar[r] & 0
  \end{tikzcd}
\end{equation*}
This has the same top row as \eqref{eq:compare-Vr}, and in fact, since we
have a map $V(r) \to \lnot D(r)$, that diagram factors through this one:
\begin{equation*}
  \begin{tikzcd}
    0 \ar[r] & R_{D(r)} := \sum_{x : R} (x = 0) \lor (r \neq 0) \ar[r,hook] \ar[d,hook]
    & R \ar[r,->>,"q"] \ar[d,equal] & R/R_{D(r)} \ar[r] \ar[d,->>] & 0 \\
    0 \ar[r] & K^{\lnot D(r)} \ar[r,hook] \ar[d,hook]
    & R \ar[r,->>,"q'"] \ar[d,equal] & R/K^{\lnot D(r)} \ar[r] \ar[d,->>] & 0 \\
    0 \ar[r] & K^{V(r)} := (r) \ar[r,hook]
    & R \ar[r,->>,"c"] & R^{V(r)} = R/(r) \ar[r] & 0
  \end{tikzcd}
\end{equation*}
\textbf{It would be nice to understand the middle row better.}
Thomas suggested that $R^{\lnot D(r)} = \lim_k R/\langle r^k \rangle$,
using the theory of infinitesimal neighbourhoods, but I'm not sure it applies here.
If true, that $K^{\lnot D(r)} = \cap_k \langle r^k \rangle$.
Other questions:
Is the map $c : R \to R^{\lnot D(r)}$ epic in this case?
Is $c(r)$ nilpotent?

Altogether there are approximately 15 modules above, five for each of
$V(r)$, $D(r)$ and $\lnot D(r)$.
(A couple turn out to be equal, and a couple were introduced that don't
come from the general construction.)
One can also take kernels or cokernels of the various vertical maps.

\textbf{Can we show that some of these are not Adelman, or have other bad properties?}
(Presumably by proving statements like: $\lnot (\text{for all $r : R$, $R_{D(r)}$ is Adelman})$.)
The technique in Theorem 3.2.1 of the topology draft might be helpful.
Note that the map $\varphi_r$ that appears above also appears there.

Methods for showing that an $R$-module $M$ is not Adelman:
\begin{itemize}
\item Show that the map $M \to M^{\ast\ast}$ is not an equivalence.
\item Show that $M$ is not wqc.
\item Show that there exists an epi $e : C \to M$ and a map $f : A \to M$,
with $A$ and $C$ finitely co-presented, such that $f$ does not lift through $e$.
(See \Cref{maps-into-coherent-lift}.)
\item Show that there exists a mono $i : M \to C$ and a map $f : M \to A$,
with $A$ and $C$ finitely presented, such that $f$ does not extend over $i$.
(The dual of \Cref{maps-into-coherent-lift}.)
\end{itemize}

When $V$ is a closed proposition, such as $V(r)$, we know that
$R^V$ is Adelman, which implies that $K^V$ and $R/K^V$ are as well.
\rednote{Citation?}

\subsection{Questions (and some answers)}
\begin{enumerate}
\item Is dualization an antiequivalence of Adelman modules?

      Yes.
\item Are Adelman modules closed under extension?

      If the quotient module is also finitely generated, then yes.
\item\label{ad-hom-closed} Are Adelman modules closed under $\Hom$?

      Yes, \cref{hom-closure}.
\item Are Adelman modules closed under $\_\otimes\_$?

      Yes, for finite free \(B\), by \cref{ad-hom-closed} and \(A\otimes B^\ast=\Hom(B,A)\).
\item Does exponentiating with closed propositions preserve Adelman?

      Yes, by partial answer to \cref{finite-exp-closed}.
\item\label{finite-exp-closed} Are projective/proper/finite-scheme products of Adelman modules Adelman?

      Partial answer:
      \begin{enumerate}[(i)]
      \item Let $X$ be a finite scheme, i.e. $X$ an affine scheme such that $R^X$ is a finitely presented $R$-module.
      For any Adelman-module $M:X\to\Mod{R}_{\mathrm{ad}}$ on $X$, $(x:X)\to M_x$ is Adelman: We can exponentiate the diagram witnessing that it is an Adelman-module and we can globalize the presentations of pointwise fp modules.
      Exponentiating is left exact and since $M$ is wqc and $X$ affine, it is also right exact.
      \end{enumerate}

\item Are the cohomology groups of Adelman module bundles on projective schemes Adelman?

      None so far. It seems a bit much to ask, but it would be really great for Serre-Duality, so it should be worth it to try to adapt e.g.\ \cite[19.1.3]{vakil}.

\item Is any Adelman-module strongly quasi-coherent?

      No. The kernel $K$ of $2\cdot\_:R\to R$ can not be strongly quasi-coherent in general.
      So far there is only an external argument: \url{https://www.youtube.com/watch?v=mdjq91Gc6Lo} around 1:10:00 states that strongly quasi-coherent modules are invariant under pullback. But in general $K$ can change between trivial and non-trivial on different stages.
\item Does being Adelman depend only on the underlying type?  Maybe related to being an fppf sheaf?
\item What are examples of non-Adelman modules?

      $R^\mathbb{N}$.  A direct sum of countable copies of $R$.  (Citations/proofs?)

Are there some that are ``smaller'' in some sense?
What about the ones mentioned in \cref{strange-modules}?
\end{enumerate}
